\chapter[奇点消解：围绕 Hironaka 定理]{奇点消解：\protect\\围绕 Hironaka 定理}\label{chap:ag-4}
% 原PDF p67.
本章在代数闭域 $k$ 上工作. $k$ 上的``簇''是有限型、不可约且既约的 $k$-概形. 记 $X^{\mathrm{sing}}$ 为 $X$ 中非正则点构成的集合，$X^{\mathrm{reg}}$ 为其补集. 已知 $X^{\mathrm{reg}}$ 是非空开集，即``光滑点开集''. 当 $X$ 正规时，已知 $X^{\mathrm{sing}}$ 是余维至少为 $2$ 的闭集.
\section{奇点消解问题}\label{sec:ag-4-1}
簇 $X$ 的一个\newterm{ag-resolution-singularities}{奇点消解}{resolution of singularities}是一个真双有理态射 $f:X'\longrightarrow X$，其中 $X'$ 是光滑簇. 一个简单的例子（却远不足以代表此问题的复杂性）是 $X$ 为曲线的情形. 此时存在唯一的奇点消解，由曲线的正规化给出. 特别地，此时 $f$ 是有限态射. 一般而言，正规化当然是通往可能的奇点消解的一项成本很低的初步操作，但远远不够.
奇点消解本身的存在性已颇有意义；但为了应用，甚至为了证明其存在性，还希望加上若干要求. 例如，可以先要求 $f$ 为射影态射.
\begin{numpara}[例外轨迹的形状]\label{item:ag-4-1-1}
对任意双有理态射 $f:X'\longrightarrow X$，存在 $X$ 中最大的非空开集 $U$，使得 $f:f^{-1}(U)\xrightarrow\sim U$. $X'$ 中 $f^{-1}(U)$ 的补集称为 $f$ 的\emph{例外轨迹}，记为 $\operatorname{Ex}(f)$. 因而它是 $X'$ 的闭集.
对``良好''的奇点消解首先可以要求：其例外轨迹尽可能简单，例如是一个除子的支集. 但一般的除子仍过于复杂（$\mathbb P^n$ 中的超曲面就可能有十分棘手的奇点）. 因此要用到\newterm{ag-normal-crossings-divisor}{正规交叉除子}{normal crossings divisor}的概念.
回顾：在极大理想为 $\mathfrak m$ 的正则局部环 $A$ 中，$\mathfrak m$ 的一族元素 $(f_1,\ldots,f_n)$，若其在 $\mathfrak m/\mathfrak m^2$ 中的像构成 $A/\mathfrak m$ 上的一组基，就称为一组\emph{正则参数系}. 当 $A=\mathcal O_{X,x}$ 时，也可以称它为 $X$ 在 $x$ 处的一组局部坐标.
\end{numpara}
\begin{definition}\label{item:ag-4-1-2}
设 $X$ 为正则概形.
\begin{itemize}
\item 有效除子 $D$ 称为\emph{正规交叉除子}，如果对每个 $x\in X$，存在局部坐标 $f_1,\ldots,f_n$ 和整数 $c_1,\ldots,c_n$，使得 $D|_{\mathcal O_{X,x}}=(f_1^{c_1}f_2^{c_2}\cdots f_n^{c_n})$.
\item 闭子概形 $Y$ 称为\emph{正规交叉除子}，如果它具有上述 $D$ 所对应的 $Y(D)$ 的形式.
\end{itemize}
\end{definition}
\begin{lemma}\label{item:ag-4-1-3}
设 $X$ 正则，$Y\subset X$ 是闭集. 下列条件等价：
\begin{enumerate}[label=\roman*)]
\item $Y$ 是上述意义下某个正规交叉除子的支集.
\item 以 $Y$ 为底集的 $X$ 的既约闭子概形，在局部由 $f_1f_2\cdots f_m=0$ 这类方程给出，其中 $f_1,\ldots,f_n$ 是 $x$ 处的局部坐标.
\item 设 $Y=\bigcup_{i\in I}Y_i$ 为 $Y$ 的不可约分支分解. 对 $I$ 的任一子集 $J$，交集 $Y_J=\bigcap_{j\in J}Y_j$ 或为空，或为余维 $|J|$ 的光滑簇.
\end{enumerate}
\end{lemma}
\begin{proof}
i) 与 ii) 的等价性显然. 假设 iii). 设 $x\in X$，并令 $J_x$ 为包含 $x$ 的分支构成的集合. 对每个 $j\in J_x$，子簇 $Y_j$ 光滑且余维为 $1$，因而局部由一个正则参数 $f_j\in\mathfrak m_x$ 给出. $Y$ 在 $x$ 处的方程于是为 $\prod_{j\in J_x}f_j$，而交集 $\bigcap_{j\in J_x}Y_j$ 由理想 $(f_j)_{j\in J_x}$ 给出. 由~\ref{item:ag-3-1-10} ii) 的证明，交集光滑这一假设保证诸 $f_j$ 在 $\mathfrak m_x/\mathfrak m_x^2$ 中线性无关. 现在假设 ii). 取 $I$ 的子集 $J$，使 $Y_J$ 非空，并取 $x\in Y_J$. 仍如上记 $J_x$，则 $J_x\supset J$. $Y$ 在 $x$ 处的方程的形式蕴含存在双射 $\{1,\ldots,m\}\xrightarrow\sim J_x$，使得对每个 $i=1,\ldots,m$，$f_{j(i)}$ 是 $Y_{j(i)}$ 的局部方程. \ref{item:ag-3-1-10} ii) 再次说明 $Y_J$ 在 $x$ 处光滑.
\end{proof}
现在对奇点消解 $f:X'\longrightarrow X$ 引入如下性质：
\begin{quote}
\textup{(RS1)} 例外轨迹 $\operatorname{Ex}(f)$ 是一个正规交叉除子的支集.
\end{quote}
对良好奇点消解的另一项自然要求是它不改变 $X$ 的正则轨迹：
\begin{quote}
\textup{(RS2)} 例外轨迹 $\operatorname{Ex}(f)$ 包含于（因而等于）$f^{-1}(X^{\mathrm{sing}})$.
\end{quote}
注意，在有 (RS2) 的情况下，还可以加强 (RS1)，要求赋予既约闭子概形结构的闭集 $X^{\mathrm{sing}}$ 的概形论原像为一个除子.
% 原PDF p68.
\begin{example}\label{item:ag-4-1-4}
设 $S\subset\mathbb A_k^3$ 为方程 $xy-z^2=0$ 所定义的二次锥. 它在 $(0,0,0)$ 处有一个孤立奇点. 在该点的爆破 $\pi:\widetilde S\longrightarrow S$ 是满足 (RS2) 与加强的 (RS1) 的奇点消解. 例外除子是一个 $\mathbb P^1$.
事实上，以三个仿射图
\begin{align*}
U_z&=\operatorname{Spec}\bigl(k[\tfrac{u}{w},\tfrac{v}{w},z]\bigr),\\
U_y&=\operatorname{Spec}\bigl(k[\tfrac{u}{v},y,\tfrac{w}{v}]\bigr),\\
U_x&=\operatorname{Spec}\bigl(k[x,\tfrac{v}{u},\tfrac{v}{w}]\bigr)
\end{align*}
覆盖 $\mathbb A^3$ 在 $(0,0,0)$ 处的爆破，可见 $S$ 的全变换依次由方程
\begin{align*}
z^2\bigl(\tfrac{u}{w}\cdot\tfrac{v}{w}-1\bigr),\qquad
y^2\bigl(\tfrac{w^2}{v^2}-\tfrac{u}{v}\bigr),\qquad
x^2\bigl(\tfrac{w^2}{u^2}-\tfrac{v}{u}\bigr)
\end{align*}
给出；除去 $z$、$y$ 或 $x$-挠部分，就得到严格变换，其方程依次为
\begin{align*}
\tfrac{u}{w}\cdot\tfrac{v}{w}-1,\qquad
\tfrac{w^2}{v^2}-\tfrac{u}{v},\qquad
\tfrac{w^2}{u^2}-\tfrac{v}{u}.
\end{align*}
这些方程表明，$U_z\cap\widetilde S$ 等于 $(U_y\cap\widetilde S)\simeq\operatorname{Spec}(k[y,\tfrac{w}{v}])=\mathbb A_k^2$ 与 $(U_x\cap\widetilde S)=\operatorname{Spec}(k[x,\tfrac{w}{u}])=\mathbb A_k^2$ 的交集.
\end{example}
更一般地，(RS1) 与 (RS2) 两条性质有助于研究孤立奇点，因为它们使得我们可以为这样的奇点关联一个组合不变量，例如例外正规交叉除子的对偶单纯复形的同伦型.
不过，其他潜在应用还需要增加新的要求.
% 原PDF p69.
\begin{numpara}[与紧化的相容性]\label{item:ag-4-1-5}
$X$ 的紧化是包含 $X$ 为开集的真 $k$-簇 $\overline X$. 如果 $X$ 拟射影，可以将 $X$ 嵌入某个 $\mathbb P_k^n$ 并取闭包而得到紧化. 然而，即使 $X$ 光滑，自然得到的紧化也很少光滑. 此外，即使 $\overline X$ 光滑，补集 $\overline X\setminus X$ 也很少光滑，而且经双有理态射后几乎不可能变得光滑. 不过仍可希望它尽可能简单. 对应用，尤其是上同调方面的应用，所寻求的性质是它成为正规交叉除子的支集.
在这种框架下，初始资料不再只是一个簇 $X$，而是一对 $(X,Z)$，其中 $Z\subsetneq X$ 是闭子集. 所希望的性质于是变为
\begin{quote}
\textup{(RS1)} 闭集 $\operatorname{Ex}(f)\cup f^{-1}(Z)$ 是一个正规交叉除子的支集.

\textup{(RS2)} 例外轨迹 $\operatorname{Ex}(f)$ 包含于 $f^{-1}(X^{\mathrm{sing}}\cup Z)$.
\end{quote}
还可进一步加强 (RS1)，要求闭子概形 $f^{-1}(Z)$ 为一个除子（此时假设 $Z$ 非空，并赋予它既约结构）. 在有 (RS2) 的情况下，也同样可以要求闭子概形 $f^{-1}(X^{\mathrm{sing}})$ 为一个除子. 注意，取 $Z=\varnothing$ 便回到先前的情形.
\end{numpara}
\begin{theorem}[Hironaka]\label{item:ag-4-1-6}
假设 $k$ 的特征为 $0$. 设 $(X,Z)$ 如上. 则存在满足 (RS1) 和 (RS2) 的奇点消解 $f:X'\longrightarrow X$.
\end{theorem}
这一版本的消解已极为有用. 但人们还可能期望存在一种``函子性''的消解构造. 下面的例子表明，要求不可过高.
\begin{example}\label{item:ag-4-1-7}
设 $S$ 为二次锥，$\widetilde S$ 为前一例中的消解. 鉴于要求 iii)，每个强消解都经由爆破 $\widetilde S$ 分解. 现设 $\varphi:\mathbb A_k^2\longrightarrow S$ 由 $(x,y)\mapsto(x^2,y^2,xy)$ 给出. $\mathbb A^2$ 的唯一强消解是恒等态射. 然而，$\varphi$ 无法提升为到 $\widetilde S$ 的态射，因为 $\varphi^{-1}(0,0,0)=(0,0)$.
\end{example}
\begin{theorem}[Hironaka ?]\label{item:ag-4-1-8}
假设 $k$ 的特征为 $0$. 存在一个对上述各对 $(X,Z)$ 作奇点消解的过程
\begin{align*}
(X,Z)\longmapsto (X'\xrightarrow{f_{X,Z}}X),
\end{align*}
该过程对光滑态射具有函子性.
\end{theorem}
一个重要应用是：若群 $G$ 作用于 $X$ 且保持闭集 $Z$，则 $(X,Z)$ 存在 $G$-等变消解，即存在一个消解 $f:X'\longrightarrow X$，其中 $X'$ 带有 $G$ 作用，而 $f$ 等变.
\begin{numpara}[转换思路：嵌入消解]\label{item:ag-4-1-9}
如前所见，代数簇间的每个射影双有理态射都是一个爆破. 然而，任意爆破的计算可能相当困难. 因此可以尝试通过相继的``简单''爆破构造奇点消解. 问题在于，真正``简单''的爆破只有光滑簇中光滑子簇的爆破. 但我们起初面对的恰恰是一个不光滑的簇……
于是考虑消解的嵌入版本. 设 $Y$ 是光滑簇 $X$ 的子簇. $Y$ 的一个\emph{嵌入消解}是射影双有理态射 $f:X'\longrightarrow X$，使得 $Y$ 的严格变换 $Y'$ 是 $Y$ 的一个奇点消解.
为了寻找这样的嵌入消解，尝试只允许光滑爆破：
\begin{quote}
\textup{(RP3)} $f$ 是沿光滑中心 $Z_j\subset X_j$ 的爆破的复合
\begin{align*}
f:X'=X_n\xrightarrow{f_{n-1}}X_{n-1}\xrightarrow{f_{n-2}}\cdots\xrightarrow{f_0}X_0=X.
\end{align*}
\end{quote}
但若想在每一步爆破中控制例外轨迹，最好要求 $Z_j$ 与 $\operatorname{Ex}(f_0\circ\cdots\circ f_{j-1})$ 的交集也尽可能简单. 由此引出下列概念：
\end{numpara}
% 原PDF p70.
\begin{definition}\label{item:ag-4-1-10}
设 $X$ 为光滑簇，$Z$ 为子簇，$D$ 为正规交叉除子. 若对每个 $x$，存在 $x$ 处的局部坐标 $f_1,\ldots,f_n$，使得 $D$ 由诸 $f_i$ 的一个单项式给出，而 $Z$ 由理想 $(f_1,\ldots,f_m)$ 给出，则称 $Z$ 与 $D$ \emph{正规交叉}.
\end{definition}
特别地，上述定义中的子簇 $Z$ 是光滑的. 此外，若 $Z$ 不包含于 $D$，则 $D\cap Z$ 是 $Z$ 上的正规交叉除子. 再者，在爆破 $\widetilde X_Z$ 中，$D$ 的全变换 $\pi_Z^{-1}(D)$ 是正规交叉除子.
\begin{quote}
\textup{(RP3')} $f$ 是沿光滑中心 $Z_j\subset X_j$ 的爆破的复合
\begin{align*}
f:X'=X_n\xrightarrow{f_{n-1}}X_{n-1}\xrightarrow{f_{n-2}}\cdots\xrightarrow{f_0}X_0=X,
\end{align*}
且 $Z_j$ 与 $\operatorname{Ex}(f_0\circ\cdots\circ f_{j-1})$ 正规交叉（注意，由归纳知这一例外轨迹确为正规交叉除子）.
\end{quote}
此时 $\operatorname{Ex}(f)$ 是一个正规交叉除子的支集. 现在，如果还希望消解 $Y'\longrightarrow Y$ 满足 (RS1)，就加上
\begin{quote}
\textup{(RP1)} 严格变换 $Y'$ 与正规交叉除子 $\operatorname{Ex}(f)$ 只有正规交叉.
\end{quote}
如果还希望 $Y'\longrightarrow Y$ 满足 (RS2)，就加上
\begin{quote}
\textup{(RP2)} $f$ 在 $Y^{\mathrm{reg}}$ 上方是同构.
\end{quote}
也可将这一要求并入 (RP3)：要求爆破 $f_j$ 的中心 $Z_j$ 与 $Y^{\mathrm{reg}}$ 不相交（若注意到此条件通过归纳蕴含 $f_0\circ\cdots\circ f_{j-1}$ 在 $Y^{\mathrm{reg}}$ 上方是同构，它便有意义）.
最后还面临粘合问题. 事实上，并非所有簇都能嵌入一个光滑簇（但所有拟射影簇当然可以）. 为了能够粘合，需要嵌入消解具有一定的函子性.
\begin{theorem}[Hironaka]\label{item:ag-4-1-11}
（$k$ 仍为特征 $0$ 的域.）存在一个过程，对上述每一对 $(X,Y)$，给出满足 (RP1)、(RP2) 和 (RP3') 的嵌入消解 $f$，而且它对光滑态射 $X'\longrightarrow X$ 与闭浸入 $X\hookrightarrow X'$ 具有函子性.
\end{theorem}
% 原PDF p71.
这一定理蕴含强奇点消解的存在性，并且该消解对光滑态射具有函子性.
\begin{numpara}[转换思路：理想的主化]\label{item:ag-4-1-12}
嵌入消解定理事实上是从关于一对 $(X,\mathcal I)$ 的定理推出的，其中 $X$ 光滑，$\mathcal I$ 为任意理想层. 原问题的转移和加强对于证明所依赖的归纳推理至关重要.
设 $(X,\mathcal I)$ 是这样的一对. 光滑 $X'$ 到 $X$ 的射影双有理态射 $f:X'\longrightarrow X$ 称为
\begin{itemize}
\item $\mathcal I$ 的\emph{主化}，如果 $f^{-1}(\mathcal I)\mathcal O_{X'}$ 是可逆理想（因而是一个有效除子）；
\item $\mathcal I$ 的\emph{单项式化}，如果 $f^{-1}(\mathcal I)\mathcal O_{X'}$ 是正规交叉除子.
\end{itemize}
若承认奇点消解的存在性，先爆破 $\mathcal I$，再消解爆破 $\widetilde X_{\mathcal I}$，就得到 $(X,\mathcal I)$ 的主化；再消解 $(\widetilde X_{\mathcal I},Z)$ 这对对象，其中 $Z$ 是例外除子的支集，便得到单项式化. 但我们反而想从主化的存在性推出奇点消解的存在性.
\end{numpara}
\begin{theorem}[Hironaka]\label{item:ag-4-1-13}
（$k$ 的特征为 $0$.）存在一个过程，对上述每一对 $(X,\mathcal I)$，给出 $\mathcal I$ 的一个单项式化 $f:X'\longrightarrow X$；其形式如 (RP3') 所规定，例外轨迹 $\operatorname{Ex}(f)$ 包含于正规交叉除子 $f^{-1}(\mathcal I)\mathcal O_{X'}$ 的支集，而且该过程对光滑态射 $Y\longrightarrow X$ 和闭浸入 $X\hookrightarrow Y$ 具有函子性.
\end{theorem}
这一定理蕴含前面两个定理. 然而值得注意的是，此处已完全没有谈到奇点消解！事实上，主化的形式（公理 (RP3')）起着首要作用.
简述它如何蕴含嵌入消解. 从光滑 $X$ 中的子簇 $Y$ 出发，记 $\mathcal I$ 为其理想. 设 $f:X'\longrightarrow X$ 是上述定理中的 $\mathcal I$ 的主化. 因而 $f$ 可分解为一系列沿光滑中心 $Z_j\subset X_j$ 的爆破
\begin{align*}
f:X'=X_n\xrightarrow{f_{n-1}}X_{n-1}\xrightarrow{f_{n-2}}\cdots\xrightarrow{f_0}X_0=X,
\end{align*}
其中 $Z_j$ 与 $\operatorname{Ex}(f_0\circ\cdots\circ f_{j-1})$ 正规交叉. 先假设 $Y$ 的余维至少为 $2$. 此时，$f$ 不可能在 $Y$ 的泛点 $\eta_Y$ 上方是同构，因此 $\operatorname{Ex}(f)$ 包含 $f^{-1}(Y)$. 于是存在最小的指标 $j$，使得 $Z_j$ 包含 $\eta:=(f_0\circ\cdots\circ f_{j-1})^{-1}(\eta_Y)$. 又因为 $\operatorname{Ex}(f)$ 包含于 $f^{-1}(Y)$，点 $\eta$ 是 $Z_j$ 的泛点. 故 $Z_j$ 是 $Y$ 经 $f_0\circ\cdots\circ f_{j-1}$ 的严格变换. 由于 $Z_j$ 是爆破中心，它必定光滑，并且是 $X$ 的一个奇点消解. 由此可知，态射 $f_0\circ\cdots\circ f_{j-1}:X_j\longrightarrow X$ 是 $X$ 的一个嵌入消解，满足 (RP1)、(RP2) 和 (RP3'). 若 $Y$ 的余维为 $1$，因而是一个除子，则 $Y$ 经 $f$ 的严格变换是 $\operatorname{Ex}(f)$ 的一个不可约分支，所以它光滑，因为 $\operatorname{Ex}(f)$ 是正规交叉除子的支集.
一旦有了函子性的嵌入消解，便有簇的函子性``抽象''消解. 为得到对 $(X,Z)$ 的消解，只须再应用一次主化.
% 原PDF p72.
\enlargethispage{3\baselineskip}
\begin{numpara}[如果特征为正]\label{item:ag-4-1-14}
在低维和某些特殊情形下，奇点消解已为人所知. 最一般的结果由 De Jong 给出. 他考虑的是变更，而不是双有理态射.
\emph{变更}是真、满且泛有限的态射.
\end{numpara}
\begin{theorem}[De Jong]\label{item:ag-4-1-15}
设 $X$ 是代数闭域上的簇，$Z\subset X$ 是闭集. 则存在一个泛平展的变更 $f:X'\longrightarrow X$，使得 $X'$ 光滑，且 $f^{-1}(Z)$ 是一个正规交叉除子的支集.
\end{theorem}
虽然没有一般的函子性性质，但在群作用存在时，此结果有等变版本. 这一定理有许多重要的上同调后果.
